\begin{enumerate}

\item Calcular las reacciones del caso real mediante equilibrio y, en su caso, compatibilidad de flechas.

\item Escribir la ley flectora real M(x) en cada intervalo.

\item Introducir una fuerza virtual unidad en el punto y sentido del desplazamiento buscado.

\item Calcular la ley flectora virtual y sumar las integrales de M por M virtual dividido entre EI para todos los tramos.

\item Verificar el resultado mediante la elástica calculada con constantes de integración y condiciones de apoyo.

\item La posición de la carga no está en el centro: dividir las integrales en los dos tramos; contrastar con el resultado de Castigliano P a\textsuperscript{2} b\textsuperscript{2}/(3 L E I).

\end{enumerate}

\noindent\textbf{Operaciones y resultados intermedios.}

\begin{enumerate}

\item Integral M por M virtual dividida por EI, tramo 1: 0 a 2000 mm.

\[\fitmath{y_{1}=\frac{5.76\times10^{12}}{3.585\times10^{12}}\approx 1{,}606695\,\mathrm{mm}}\]

\item Integral M por M virtual dividida por EI, tramo 2: 2000 a 5000 mm.

\[\fitmath{y_{2}=\frac{8.64\times10^{12}}{3.585\times10^{12}}\approx 2{,}410042\,\mathrm{mm}}\]

\item Desplazamiento total por suma de tramos.

\[\fitmath{y_{3}=\frac{5.76\times10^{12}}{3.585\times10^{12}}+\frac{8.64\times10^{12}}{3.585\times10^{12}}\approx 4{,}016736\,\mathrm{mm}}\]

\end{enumerate}
